\chapter{Metodologia} \label{cap:proposta}

    A metodologia adotada neste trabalho estrutura-se de forma progressiva, articulando modelagem matemática, implementação computacional e análise de desempenho com vistas ao refinamento do formalismo proposto. Essa progressão organiza-se em duas etapas: uma primeira fase de verificação de viabilidade, já concluída e com resultados preliminares obtidos, dedicada à construção e validação dos Hamiltonianos propostos; e uma segunda fase, em desenvolvimento, voltada à implementação de algoritmos de otimização quântica variacional sobre os formalismos estabelecidos.

    Inicialmente, o trabalho concentrou-se na construção de Hamiltonianos capazes de representar problemas de otimização combinatória em um contexto de computação quântica de variáveis contínuas. Essa etapa teve como propósito estabelecer uma correspondência entre a estrutura do problema e sua codificação em operadores compatíveis com sistemas quânticos de variáveis contínuas.

    O TSP foi adotado, nesse contexto, como problema base para validação, por constituir uma instância clássica de otimização combinatória com estrutura suficientemente rica para testar o método \cite{hartmanis1982computers}, mas com complexidade ainda tratável do ponto de vista analítico e computacional. Como ponto de partida, foi considerada uma primeira abordagem baseada em uma correspondência direta entre a modelagem discreta usual e o formalismo de variáveis contínuas. Nesse contexto, cada qubit, originalmente associado a uma variável binária no modelo clássico, foi mapeado para um qumode, de modo a reproduzir o comportamento discreto esperado \cite{verdon2019quantum}.

    Para isso, foi construído um Hamiltoniano análogo ao utilizado em formulações do tipo QUBO em qubits \cite{lucas2014ising}. Considera-se um conjunto de $N$ cidades em um grafo, com $d_{ij}$ representando a distância entre as cidades $i$ e $j$, e define-se a variável binária

    $$
    x_{i,k} =
    \begin{cases}
    1, & \text{se a cidade $i$ é visitada na posição $k$},\\
    0, & \text{caso contrário}.
    \end{cases}
    $$

    A função objetivo do TSP pode então ser escrita como:

    $$
        \label{eq:objective_function_hamiltonian}
        H = \sum_{k=1}^N \sum_{i=1}^N \sum_{j=1}^N d_{ij}\, x_{i,k} x_{j,k+1}
    $$

    Sujeito às restrições:

    \begin{enumerate}
        \item Cada posição deve conter exatamente uma cidade:
        $$
            \sum_{i=1}^N x_{i,k} = 1, \qquad \forall k
        $$

        \item Cada cidade deve ser visitada exatamente uma vez:
        $$
            \sum_{k=1}^N x_{i,k} = 1, \qquad \forall i
        $$
    \end{enumerate}

    No contexto de variáveis contínuas, as variáveis discretas $x_{i,k}$ são substituídas por operadores $\hat{x}_{i,k}$ associados aos qumodes. Introduzem-se, adicionalmente, os operadores canônicos de posição $\hat{x}_{i,k}$ e momento $\hat{p}_{i,k}$, satisfazendo as relações de comutação \cite{weedbrook2012gaussian, verdon2019quantum}:

    $$
    [\hat{x}_{i,k}, \hat{p}_{j,l}] = i\hbar\,\delta_{ij}\delta_{kl}.
    $$
    \noindent adota-se, para simplificação algébrica, a convenção $\hbar=1$. Uma vez que o operador posição possui espectro contínuo, impõe-se um comportamento efetivamente binário por meio da transformação \cite{arrazola2019machine}

    $$
    \hat{x}_{i,k} = \frac{1 + \hat{x}'_{i,k}}{2},
    $$

    \noindent de modo que $\hat{x}'_{i,k} \approx \pm 1$ corresponde a $\hat{x}_{i,k} \in \{0,1\}$. Essa restrição é imposta energeticamente através de um potencial de poço duplo:

    $$
    H_{\text{bin}} = \sum_{i,k} \left( (\hat{x}'_{i,k})^2 - 1 \right)^2.
    $$

    Além disso, considera-se um termo cinético \cite{verdon2019quantum} dado por:

    $$
    \sum_{i,k} \frac{\hat{p}_{i,k}^2}{2m}.
    $$

    Dessa forma, o Hamiltoniano completo do sistema é definido como:

    \begin{equation}
    \begin{aligned}
    H &= \sum_{k=1}^N \sum_{i=1}^N \sum_{j=1}^N d_{ij}\, \hat{x}_{i,k}\hat{x}_{j,k+1}
        + \alpha \sum_{i,k} \left( (\hat{x}'_{i,k})^2 - 1 \right)^2 \\
        &\quad + \beta \sum_{i=1}^N \left(\sum_{k=1}^N \hat{x}_{i,k} - 1 \right)^2
        + \gamma \sum_{k=1}^N \left(\sum_{i=1}^N \hat{x}_{i,k} - 1 \right)^2
        + \sum_{i,k} \frac{\hat{p}_{i,k}^2}{2m},
    \end{aligned}
    \end{equation}

    \noindent onde $\alpha$, $\beta$, $\gamma$ e $m$ são hiperparâmetros do modelo \cite{lucas2014ising}. A solução do problema é obtida pela minimização do valor esperado do Hamiltoniano:

    \begin{equation}
    \min_{\psi \in \mathcal{H}} \bra{\psi} H \ket{\psi}.
    \end{equation}

    Essa formulação inicial teve como principal objetivo verificar a consistência da representação com variáveis contínuas em relação aos modelos tradicionais. Para isso, os resultados obtidos foram comparados tanto com soluções provenientes de uma busca exaustiva clássica quanto com implementações do Hamiltoniano equivalente em qubits, formulação da qual o modelo de qumodes deriva diretamente, por substituição das variáveis binárias por operadores de posição com comportamento binário imposto energeticamente. Em ambos os casos, os estados de energia mínima encontrados no modelo com qumodes coincidiram com os das abordagens de referência, indicando que a modelagem proposta preserva corretamente a estrutura do problema.

    A formulação descrita constitui a primeira etapa metodológica deste trabalho, tendo sido desenvolvida com o propósito de verificar a viabilidade da representação de problemas combinatórios no formalismo de variáveis contínuas, antes de avançar para codificações mais naturais ao paradigma de qumodes. Os resultados obtidos com essa modelagem, bem como a análise comparativa com as abordagens de referência em qubits e busca exaustiva clássica, são apresentados e discutidos no \autoref{cap:resultados}.
    
    